\section{Interval orders}\label{app:C.13}

Two materials whose figures of merit have overlapping 68\% intervals cannot be put in order, and a third may be resolved from one of them but not from the other. \Cref{sec:17.6} groups materials into tiers that respect exactly what the intervals resolve. The mathematics is that of an interval order, a partial order built from intervals, and of the heights of its elements.

\subsection{Partial orders from intervals}

A \emph{strict partial order} $\succ$ on a finite set is a relation that is irreflexive (never $a \succ a$) and transitive ($a \succ b$ and $b \succ c$ imply $a \succ c$). Together the two properties imply asymmetry: $a \succ b$ and $b \succ a$ would give $a \succ a$. Two elements with neither $a \succ b$ nor $b \succ a$ are \emph{incomparable}. A partial order differs from a ranking in that incomparability is allowed and is not transitive.

Let each element $a$ carry an interval $I_a = [\ell_a, u_a]$, oriented so that larger values are better; for a figure of merit where smaller is better, use its negative logarithm. Define
\begin{equation}
  a \succ b \quad\Longleftrightarrow\quad \ell_a > u_b ,
  \label{eq:C.13-order}
\end{equation}
that is, $a$ is resolved above $b$ when the whole interval of $a$ lies above the whole interval of $b$. The relation is irreflexive, since $\ell_a \le u_a$, and transitive: $\ell_a > u_b \ge \ell_b > u_c$ gives $\ell_a > u_c$. It is therefore a strict partial order, called an \emph{interval order}. Two elements are incomparable exactly when their intervals overlap.

Incomparability is not transitive. With $I_a = [3, 5]$, $I_b = [2, 4]$ and $I_c = [1, 2.5]$, $a$ overlaps $b$ and $b$ overlaps $c$, but $\ell_a = 3 > 2.5 = u_c$, so $a \succ c$. A ranking that declares $a$ tied with $b$ and $b$ tied with $c$ must then either tie $a$ with $c$, which discards a resolved difference, or break one of the ties, which asserts a difference the intervals do not show.

Interval orders have one property that general partial orders lack: if $a \succ b$ and $c \succ d$, then $a \succ d$ or $c \succ b$. To prove the property, suppose $a \not\succ d$, so $\ell_a \le u_d$. Then $\ell_c > u_d \ge \ell_a > u_b$, and $c \succ b$. Conversely, every finite partial order with this property can be represented by intervals as in \cref{eq:C.13-order}~\cite{fishburn1970intransitive}. When all intervals have the same length the order is called a semiorder; intervals of different widths, which arise when inputs of different provenance carry different uncertainties, give the general case.

\subsection{Heights as tiers}

A \emph{chain} above $a$ is a sequence $b_1 \succ b_2 \succ \dots \succ b_k \succ a$. The \emph{height} of $a$ is
\begin{equation}
  h(a) = 1 + \bigl(\text{length } k \text{ of the longest chain above } a\bigr) ,
  \label{eq:C.13-height}
\end{equation}
so that $h(a) = 1$ when nothing is resolved above $a$. The \emph{tier} $k$ is the set of elements of height $k$.

Each tier is an \emph{antichain}: no two of its members are comparable. If $a \succ b$, any chain above $a$ extended by $a$ is a chain above $b$, one element longer, so $h(b) \ge h(a) + 1$ and $a$, $b$ lie in different tiers. Members of one tier are therefore unresolved from each other, and calling them tied within the tier is exact. Every member $a$ of tier $k > 1$ is resolved below at least one member of tier $k - 1$. A longest chain above $a$ has $k - 1$ elements, and its last element $b$ has height exactly $k - 1$. The chain above $b$ formed by the first $k - 2$ elements shows $h(b) \ge k - 1$, while a longer chain above $b$ would extend through $b$ to a chain above $a$ of more than $k - 1$ elements.

For interval orders the heights are computed by sorting. Process the elements in decreasing order of $u_a$. When $a$ is reached, every $b$ with $b \succ a$ has $\ell_b > u_a$, so $u_b \ge \ell_b > u_a$ and $b$ has already been processed; then $h(a) = 1 + \max\{h(b) : \ell_b > u_a\}$, with the maximum over an empty set equal to zero. The cost is at most quadratic in the number of elements.

\subsection{Interpretation of tiers}

Take four materials with intervals $I_A = [8, 10]$, $I_B = [5, 9]$, $I_C = [6, 7]$ and $I_D = [1, 9.5]$. Nothing lies above $A$, $B$ or $D$, since no lower end exceeds their upper ends ($\ell_A = 8 < 9 \le u_B, u_D$, and $\ell_B, \ell_C, \ell_D < 9$), so all three have height 1. Only $A$ is resolved above $C$ ($8 > 7$), so $C$ has height 2. The tiers are $\{A, B, D\}$ and $\{C\}$. Material $D$, whose interval is the widest, is in the first tier although its midpoint, 5.25, is below that of $C$, 6.5: nothing is resolved above $D$. $C$ is in the second tier although $B$, which is in the first tier, is not resolved above $C$.

A tier does not assert that every member of a higher tier beats every member of a lower one: $B$ does not beat $C$. Nor does it assert that a member of a higher tier is likely to be better than one of a lower tier, since a wide interval rises to the first tier because it is wide. Tiers also depend on the confidence level of the intervals and on the widths assigned to the inputs; at a higher level, the intervals widen and the tiers merge.
